\chapter{Cortisol}

\section*{What Is This Test?}
Cortisol (hydrocortisone) is the primary glucocorticoid hormone produced by the zona fasciculata of the adrenal cortex. It is synthesized from cholesterol through a series of enzymatic reactions, with the rate-limiting step being the transport of cholesterol into mitochondria by steroidogenic acute regulatory protein (StAR), followed by side-chain cleavage by CYP11A1. Cortisol secretion is regulated by the hypothalamic-pituitary-adrenal (HPA) axis: corticotropin-releasing hormone (CRH) from the hypothalamus stimulates adrenocorticotropic hormone (ACTH) release from the anterior pituitary, which in turn stimulates cortisol synthesis and secretion. Cortisol exerts negative feedback on both CRH and ACTH. Cortisol circulates in blood approximately 90--95\% protein-bound: primarily to corticosteroid-binding globulin (CBG, transcortin, 75--80\%) and albumin (15\%), with only 5--10\% free and biologically active. Cortisol follows a pronounced circadian rhythm with peak levels in the early morning (6--8 AM, approximately 10--25 mcg/dL) and a nadir around midnight (<5 mcg/dL). Cortisol is essential for life, regulating glucose metabolism (promoting gluconeogenesis), suppressing inflammation, maintaining vascular tone, modulating immune function, and facilitating the stress response. Cortisol testing can be performed on serum, saliva (late-night salivary cortisol), or urine (24-hour urinary free cortisol).

\section*{Alternative Names}
Hydrocortisone; Serum Cortisol; Plasma Cortisol; 24-Hour Urinary Free Cortisol; Late-Night Salivary Cortisol; Compound F; Cortisol AM/PM; Dexamethasone Suppression Test; ACTH Stimulation Test (Cosyntropin Test)

\section*{Principle}
Serum cortisol is measured using competitive immunoassays on automated analyzers \cite{rifai2018tietz}. The method uses an anti-cortisol antibody that competes for cortisol in the patient sample with a cortisol analog labeled with a chemiluminescent tag (ruthenium on Roche ECLIA, acridinium ester on Abbott/Siemens). A displacement agent (danazol, 8-anilino-1-naphthalene-sulfonic acid) releases cortisol from CBG and albumin. The signal is inversely proportional to cortisol concentration. Common platforms: Roche Cobas ECLIA, Abbott Architect CMIA, Siemens ADVIA Centaur/Atellica CLIA, and Beckman Coulter DxI. Cortisol immunoassays have cross-reactivity with other steroids (prednisolone 30--50\%, 11-deoxycortisol 5--15\%, cortisone 0.3--3\%, prednisone 2--5\%), though modern assays have improved specificity. The gold standard reference method is isotope dilution liquid chromatography-tandem mass spectrometry (LC-MS/MS), which is more specific than immunoassay and can simultaneously measure multiple steroids. Salivary cortisol is measured by highly sensitive enzyme immunoassay (ELISA) or LC-MS/MS and reflects free cortisol levels because CBG is absent in saliva. Twenty-four-hour urinary free cortisol requires extraction and immunoassay or LC-MS/MS.

\section*{History}
Cortisol was first isolated by Edward Kendall from adrenal extracts at the Mayo Clinic in 1936 (for which he shared the 1950 Nobel Prize with Reichstein and Hench). Tadeus Reichstein independently isolated and synthesized cortisol in 1937. The first clinical use of cortisone for rheumatoid arthritis was reported by Philip Hench in 1949, demonstrating its dramatic anti-inflammatory effects. The radioimmunoassay for cortisol was developed by Murphy in 1967. The HPA axis was elucidated by Geoffrey Harris, Roger Guillemin, and Andrew Schally in the 1950s--1970s (Guillemin and Schally shared the 1977 Nobel Prize for discovering peptide hormone production in the brain). The dexamethasone suppression test was first reported by Liddle in 1960 for diagnosing Cushing syndrome. The ACTH stimulation test (cosyntropin test) was developed in the 1960s. Late-night salivary cortisol was validated as a screening test for Cushing syndrome in the 1990s. The 2008 Endocrine Society guidelines for Cushing syndrome and 2016 guidelines for adrenal insufficiency codified current diagnostic algorithms \cite{bornstein2016diagnosis}.

\section*{Why This Test Is Ordered}
\begin{itemize}
  \item Screening for Cushing syndrome in patients with central obesity, facial rounding (moon face), dorsocervical fat pad (buffalo hump), purple striae >1 cm wide, proximal muscle weakness, easy bruising, osteoporosis, and resistant hypertension --- initial screening tests: 24-hour urinary free cortisol (at least 2 measurements), late-night salivary cortisol (2 measurements), or 1 mg overnight dexamethasone suppression test \cite{nieman2015treatment}
  \item Diagnosis of adrenal insufficiency (primary or secondary) in patients with unexplained weight loss, fatigue, hypotension, hyponatremia, hyperkalemia, hyperpigmentation (primary only), or eosinophilia --- 250 mcg cosyntropin (ACTH 1-24) stimulation test preferred \cite{bancos2015diagnosis}
  \item Monitoring glucocorticoid replacement therapy in Addison disease or congenital adrenal hyperplasia (target: hydrocortisone 15--25 mg/day divided)
  \item Detection of congenital adrenal hyperplasia by elevated 17-hydroxyprogesterone and impaired cortisol response to ACTH stimulation
  \item Evaluation of critical illness-related corticosteroid insufficiency (CIRCI) in septic shock unresponsive to vasopressors --- random cortisol <15 mcg/dL or delta cortisol <9 mcg/dL after 250 mcg cosyntropin suggests CIRCI
  \item Assessing HPA axis recovery after prolonged glucocorticoid therapy (>3 weeks at prednisone-equivalent >5 mg/day)
  \item Investigation of ectopic ACTH syndrome (small cell lung cancer, bronchial carcinoid) with severe hypercortisolism, hypokalemia, and metabolic alkalosis
  \item Differentiating pseudo-Cushing states (depression, alcoholism, obesity) from true Cushing syndrome
\end{itemize}

\section*{Is This a Primary or Secondary Test}
Serum cortisol is a primary screening test for both Cushing syndrome and adrenal insufficiency, but under dynamic testing conditions (suppression or stimulation) rather than as a single random measurement. The 1 mg overnight dexamethasone suppression test is the most widely used outpatient screening test for Cushing syndrome. The cosyntropin (ACTH) stimulation test is the gold standard for diagnosing primary adrenal insufficiency. Random serum cortisol has limited diagnostic utility because of circadian variation, stress-induced elevations, and wide inter-individual variability \cite{raff2003physiologic}. The Endocrine Society recommends against using random serum cortisol for screening of adrenal disorders \cite{findling2017diagnosis}.

\section*{How This Test Is Performed}
For serum cortisol, a healthcare professional draws a blood sample from a vein into a serum separator tube (gold-top) or plain red-top tube, ideally at 8 AM (peak) and/or 4 PM or midnight (nadir). The timing MUST be documented on the requisition. For the 1 mg overnight dexamethasone suppression test, the patient takes 1 mg dexamethasone orally at 11 PM, and blood is drawn at 8 AM the next morning for serum cortisol measurement. Cortisol <1.8 mcg/dL effectively excludes Cushing syndrome (sensitivity 95\%, specificity 80\%). For the ACTH stimulation test, baseline cortisol and ACTH are drawn, cosyntropin 250 mcg is administered IV or IM, and cortisol is measured at 30 and 60 minutes. Peak cortisol >18--20 mcg/dL indicates adequate adrenal reserve. For late-night salivary cortisol, the patient collects saliva at 11 PM--midnight using a cotton swab (Salivette) and the sample is refrigerated until testing. For 24-hour urinary free cortisol, the patient collects all urine for 24 hours in a jug, which is refrigerated, and total volume and creatinine are measured. Saliva and urine testing are performed on multiple separate days.

\section*{What Preparation Is Needed}
For morning serum cortisol: no fasting required but the sample must be drawn between 6--9 AM. For dexamethasone suppression test: patient must be compliant with dexamethasone ingestion, and the sample drawn at 8 AM the next day. Drugs that accelerate dexamethasone metabolism (phenytoin, phenobarbital, carbamazepine, rifampin) cause false-positive results by reducing dexamethasone levels. Estrogen-containing contraceptives and pregnancy increase CBG, elevating total serum cortisol --- a normal suppression test in this context requires a higher cortisol threshold. For salivary cortisol: no brushing teeth, eating, drinking, or smoking for 30 minutes before collection; avoid contamination with blood from gingival disease. For the ACTH stimulation test: the patient should hold hydrocortisone or prednisone for 24 hours before testing (dexamethasone can be continued as it does not cross-react in cortisol assays). High-dose biotin (>5 mg/day) must be discontinued 48 hours before testing. Profound stress (acute illness, trauma, surgery) elevates cortisol and invalidates most HPA axis testing.

\section*{Contraindications and Precautions}
Factors elevating total serum cortisol: estrogen therapy, oral contraceptives, pregnancy (increase CBG 2--3 fold), acute stress (illness, hospitalization, pain, anxiety), severe depression (pseudo-Cushing), poorly controlled diabetes, alcoholism (especially during withdrawal), chronic kidney disease, and obesity. Dexamethasone suppression test false-positives: drugs increasing dexamethasone metabolism (phenytoin, phenobarbital, carbamazepine, rifampin, pioglitazone), malabsorption, non-compliance, estrogen therapy (higher CBG). False-negative dexamethasone tests (suppression despite Cushing): cyclical Cushing, some cases of mild Cushing disease. Cosyntropin test considerations: recent ACTH deficiency (<4--6 weeks' duration) may still produce a normal cortisol response because the adrenal glands have not yet atrophied. If secondary adrenal insufficiency is suspected (ACTH deficiency of recent onset), the insulin tolerance test or metyrapone test may be more sensitive. Interpretation of salivary cortisol is affected by shift work, irregular sleep patterns, cross-reactivity with topical hydrocortisone, and contamination with steroid-containing creams. Hemolyzed or icteric specimens interfere with spectrophotometric detection.

\section*{Risks}
Minimal risk from venipuncture. Cosyntropin (ACTH 1-24) administration: very rarely allergic reactions (cosyntropin is a synthetic peptide and does not contain animal proteins). Dexamethasone side effects from a single 1 mg dose are negligible. For the insulin tolerance test (rarely used): risk of severe hypoglycemia requiring continuous monitoring, IV glucose rescue, and physician supervision. Salivary collection: no risks.

\section*{What Happens After the Test}
Serum cortisol results are available within 1--4 hours; salivary cortisol within 1--3 days; 24-hour urinary free cortisol within 3--5 days. Abnormal screening tests require confirmation: if the 1 mg dexamethasone suppression test is abnormal (cortisol >1.8 mcg/dL), proceed to 24-hour urinary free cortisol (2 measurements), late-night salivary cortisol (2 measurements), or the 2-day low-dose dexamethasone suppression test (0.5 mg every 6 hours for 48 hours). If 2 or more first-line tests are abnormal, ACTH measurement, pituitary MRI, and possibly inferior petrosal sinus sampling (IPSS) are performed to determine the etiology of ACTH-dependent Cushing syndrome. If the ACTH stimulation test is abnormal (peak cortisol <18 mcg/dL), adrenal insufficiency is confirmed; the next step is measuring ACTH to determine primary vs. secondary etiology. Patients with confirmed adrenal insufficiency require glucocorticoid replacement (hydrocortisone 15--25 mg/day in divided doses) and stress-dose steroid education.

\section*{Understanding Results}

\subsection*{Below Normal}
\begin{itemize}
  \item \textbf{Morning cortisol <3 mcg/dL:} Strongly suggests adrenal insufficiency. If confirmed by abnormal cosyntropin stimulation test (peak cortisol <18 mcg/dL at 30 or 60 minutes), the diagnosis is adrenal insufficiency.
  \item \textbf{Primary adrenal insufficiency (Addison disease):} Low cortisol, elevated ACTH (usually >100 pg/mL, often >500 pg/mL). The most common cause in developed countries is autoimmune adrenalitis (80\%), with anti-21-hydroxylase antibodies positive in >80\% of cases. Other causes: tuberculosis (most common worldwide), adrenal hemorrhage (Waterhouse-Friderichsen syndrome in meningococcemia), adrenal metastases, congenital adrenal hyperplasia, adrenoleukodystrophy, and medications (ketoconazole, etomidate, mitotane, checkpoint inhibitors).
  \item \textbf{Secondary adrenal insufficiency:} Low cortisol with inappropriately normal or low ACTH (<20 pg/mL). Most commonly from chronic exogenous glucocorticoid use causing HPA axis suppression. Other causes: pituitary adenoma (non-functioning macroadenoma compressing normal tissue), pituitary surgery/radiation, Sheehan syndrome, traumatic brain injury, lymphocytic hypophysitis, and immune checkpoint inhibitor-related hypophysitis.
  \item \textbf{Suppressed cortisol on dexamethasone suppression test (<1.8 mcg/dL):} Normal --- effectively excludes Cushing syndrome. Also seen in patients with adrenal insufficiency.
  \item \textbf{Congenital adrenal hyperplasia:} Low cortisol with elevated ACTH and elevated steroid precursors (17-hydroxyprogesterone in 21-hydroxylase deficiency, which accounts for >90\% of CAH cases).
\end{itemize}

\subsection*{Above Normal}
\begin{itemize}
  \item \textbf{Non-suppressed cortisol on 1 mg dexamethasone suppression test (>1.8 mcg/dL):} Abnormal --- requires further testing for Cushing syndrome. Sensitivity 95\%, specificity 80\% using the 1.8 mcg/dL cutoff. If 2 first-line tests are abnormal, Cushing syndrome is confirmed.
  \item \textbf{Cushing disease (ACTH-secreting pituitary adenoma):} Approximately 70\% of endogenous Cushing syndrome. Elevated or inappropriately normal ACTH, elevated cortisol, lack of suppression on low-dose dexamethasone, but suppression (>50\%) on high-dose dexamethasone (8 mg overnight). Pituitary MRI detects adenoma in only 50--60\% of cases; inferior petrosal sinus sampling (IPSS) with CRH stimulation is the gold standard for localization (central-to-peripheral ACTH gradient >2:1 basal, >3:1 post-CRH).
  \item \textbf{Ectopic ACTH syndrome:} Approximately 10--15\% of endogenous Cushing. Markedly elevated ACTH (often >200 pg/mL), severe hypercortisolism, hypokalemia (<3.0 mmol/L), metabolic alkalosis. Most commonly: small cell lung carcinoma, bronchial carcinoid, thymic carcinoid, pancreatic neuroendocrine tumors, medullary thyroid carcinoma, pheochromocytoma. No suppression on either low- or high-dose dexamethasone.
  \item \textbf{Adrenal Cushing syndrome:} Approximately 15--20\% of endogenous Cushing. Adrenal adenoma (10\%) or carcinoma (5--8\%) autonomously secreting cortisol. Low/undetectable ACTH (<10 pg/mL), cortisol not suppressed on either low- or high-dose dexamethasone. Adrenal CT shows unilateral adenoma or heterogeneous mass; carcinomas are typically >6 cm with irregular borders.
  \item \textbf{Physiologic hypercortisolism:} Pregnancy (placental CRH stimulates HPA axis, cortisol increases 2--3 fold by third trimester), severe stress (illness, surgery, burns, trauma), strenuous exercise. Cortisol may exceed 50 mcg/dL in severe stress.
  \item \textbf{Pseudo-Cushing states:} Major depressive disorder (especially melancholic type), chronic alcoholism (active drinking or withdrawal), poorly controlled diabetes mellitus, severe obesity, anorexia nervosa \cite{lacroix2015cushing}. Cortisol circadian rhythm may be blunted, and dexamethasone suppression may be abnormal. Investigation should be deferred until the underlying condition is treated.
  \item \textbf{Elevated CBG (estrogen effect):} Oral contraceptives, estrogen replacement therapy, tamoxifen, pregnancy --- total cortisol elevated 2--3 fold. Free cortisol (salivary or urinary) is normal. Cortisol suppression on dexamethasone testing may be blunted, so a higher threshold is used (cortisol <5 mcg/dL after 1 mg dexamethasone).
\end{itemize}

\section*{Normal Results}
\begin{itemize}
  \item \textbf{Serum cortisol, 8 AM:} 5--25 mcg/dL (140--690 nmol/L) --- varies by assay; peak occurs at 6--8 AM
  \item \textbf{Serum cortisol, 4 PM:} 3--15 mcg/dL (80--410 nmol/L) --- approximately 50\% of the 8 AM value
  \item \textbf{Serum cortisol, midnight (sleeping):} <5 mcg/dL (<140 nmol/L) --- loss of circadian rhythm is an early sign of Cushing syndrome
  \item \textbf{1 mg overnight dexamethasone suppression test:} Cortisol <1.8 mcg/dL (<50 nmol/L) at 8 AM --- threshold for excluding Cushing syndrome
  \item \textbf{ACTH stimulation test (250 mcg cosyntropin IV/IM):} Baseline cortisol >5 mcg/dL; peak cortisol >18--20 mcg/dL at 30 or 60 minutes (adequate adrenal reserve)
  \item \textbf{24-hour urinary free cortisol (adults):} <50 mcg/24 hours (Roche ECLIA; lab-specific range). Values >4 times upper limit of normal are diagnostic for Cushing syndrome
  \item \textbf{Late-night salivary cortisol (11 PM):} <0.1 mcg/dL (<2.8 nmol/L; lab-specific). Two elevated values are highly sensitive for Cushing syndrome
  \item \textbf{Free cortisol (saliva or calculated free cortisol index):} Saliva reflects free cortisol accurately; approximately 10\% of total cortisol is free
  \item \textbf{Pregnancy (total cortisol):} Increases 2--3 fold above non-pregnant values from second trimester onward; use trimester-specific ranges or salivary cortisol
\end{itemize}

\section*{Follow-up Tests}
\begin{itemize}
  \item \textbf{ACTH (adrenocorticotropic hormone):} Essential for classifying adrenal insufficiency or Cushing syndrome as ACTH-dependent vs. ACTH-independent. Must be drawn simultaneously with cortisol (8 AM) into a pre-chilled EDTA tube, placed on ice, centrifuged at 4 degrees C, and frozen within 1 hour.
  \item \textbf{24-hour urinary free cortisol:} Screening test for Cushing syndrome measuring integrated cortisol excretion, reflecting cortisol secretion over 24 hours. At least 2 collections required. Elevated more than 3--4 times the upper limit of normal confirms Cushing syndrome.
  \item \textbf{Late-night salivary cortisol:} Screening test for Cushing syndrome based on loss of circadian cortisol rhythm. Two separate midnight collections. Highly sensitive and specific; convenient outpatient test.
  \item \textbf{Low-dose dexamethasone suppression test (2-day, 2 mg):} Dexamethasone 0.5 mg every 6 hours for 48 hours; serum cortisol measured at 48 hours. Cortisol <1.8 mcg/dL excludes Cushing syndrome.
  \item \textbf{High-dose dexamethasone suppression test (8 mg overnight):} Dexamethasone 8 mg at 11 PM, cortisol at 8 AM. Cortisol suppression >50\% from baseline suggests Cushing disease (pituitary); lack of suppression suggests ectopic ACTH or adrenal Cushing.
  \item \textbf{CRH stimulation test:} Ovine CRH 1 mcg/kg IV; ACTH and cortisol measured at baseline and at 15, 30, 45, 60, 90, and 120 minutes. ACTH increase >35--50\% and cortisol >20\% at 15--30 minutes suggests Cushing disease.
  \item \textbf{Inferior petrosal sinus sampling (IPSS) with CRH stimulation:} Gold standard for differentiating pituitary from ectopic ACTH source in ACTH-dependent Cushing syndrome. Central-to-peripheral ACTH gradient >2:1 basal and >3:1 after CRH confirms Cushing disease.
  \item \textbf{Pituitary MRI with gadolinium:} Detects pituitary adenomas. T1-weighted post-contrast images are most sensitive. A pituitary microadenoma (<1 cm) is identified in only 50--60\% of Cushing disease cases.
  \item \textbf{Adrenal CT (non-contrast, with washout):} For ACTH-independent Cushing syndrome. Adenomas are typically homogeneous, <4 cm, with lipid-rich attenuation <10 Hounsfield units and >50\% washout on delayed imaging. Carcinomas are >6 cm, heterogeneous, irregular.
  \item \textbf{Anti-21-hydroxylase antibodies:} Confirms autoimmune etiology in primary adrenal insufficiency. Positive in >80\% of autoimmune Addison disease.
\end{itemize}

\section*{References}
\bibliographystyle{unsrtnat}
\bibliography{references}

\clearpage
